\chapter{序説 (Introductory)}

\section{数学と生命の探求}

カントは同時代の化学について、「真の科学とは、数学が適用されうる度合いに比例して存在するにすぎない」と述べた。この厳格な基準に照らせば、化学は当時は真の科学ではなく、経験的な技術や博物学にとどまっていた。しかし化学が物理学や数学の原理と融合した今日、化学は真の精密科学へと昇華した。

生命の形態を研究する生物学、すなわち形態学（Morphology）においても、同様の転換が要請されている。古生物学や比較解剖学は、生命の形態の膨大な記述や系統発生の追跡において偉大な業績を上げてきた。だが、われわれが「なぜその形態をとるのか」「いかなる物理的・力学的必然性によってその輪郭が生み出されたのか」を問うとき、古典的な記載生物学の語彙だけでは不十分である。

生物の形態とは、成長の途上において物質に作用した力の痕跡、すなわち「力学的ダイアグラム（Diagram of Forces）」である。力学の言葉を借りずして形態の本質を記述することはできず、数学の言葉を用いずして力の作用を分析することはできない。

\section{力と形態の関係}

物理学において、力（Force）とは運動の変化をもたらす作用である。静止した物質や平衡状態にある構造において、力は内的な応力や表面張力、弾性抵抗として物質の配列と幾何学的配置を決定づける。

生命体であっても物質である以上、外部から作用する重力、大気圧、水圧、流体の摩擦、さらには分子間に働く表面張力、毛細管現象、浸透圧、電磁気力といった普遍的な物理法則の支配を完全に受けている。細胞分裂における紡錘体の配向、細胞膜の接触角、石灰質骨片の結晶成長、骨にかかる主応力線の軌跡に至るまで、生命のあらゆる造形美は、物質と力との相互作用の精妙な幾何学的帰結である。

われわれの課題は、自然淘汰という究極要因（進化の歴史）を性急に持ち出す前に、まず至近要因としての力学的・幾何学的制約を明らかにすることである。物理法則によってそもそも存在し得ない形態は淘汰の対象にすらならず、物理法則が必然的に導く形態は、適応の有無にかかわらず自ずと実現されるからである。

\section{本書の構成と展望}

本書において、われわれはまず**「大きさ（Magnitude）」**の物理的効果から出発する。長さ、面積、体積の幾何学的スケーリングが、重力と表面張力の勢力関係をいかに一変させるかを論じる。

続いて、成長の速度と時間的ダイナミクス、単一細胞の形態と表面張力の関係、多数の細胞が凝集した組織における幾何学的充填問題（プラトーの境界条件と120度接合）、放散虫や海綿の美しい骨片骨格、オウムガイやアンモナイトに見られる対数螺旋、植物の葉序（フィボナッチ幾何学）、そして脊椎動物の骨格構造における力学的最適性（主応力線と片持ち梁の原理）を順次検討する。

そして最後に、直交座標系の平滑な変形によって、ある生物の形態から近縁種の形態へと幾何学的に写像しうる**「変形の理論（Theory of Transformations）」**を展開する。幾何学と物理法則という普遍のレンズを通して、生命の多様な形態の中に潜む深遠な統一性を描き出すことが、本書の目指す地平である。
